\documentclass[../../main.tex]{subfiles}
\begin{document}
\section{切--余切同构}\label{section:ch13-tangent-cotangent}

每个 Riemann 度量的另一个便利特点是：它在切丛与余切丛之间提供了一个自然同构。给定带边或不带边光滑流形 $M$ 上的一个 Riemann 度量 $g$，我们定义一个丛同态 $\hat g:TM\to T^*M$ 如下。对每个 $p\in M$、每个 $v\in T_pM$，令 $\hat g(v)\in T_p^*M$ 为由
\begin{align*}
\hat g(v)(w)=g_p(v,w)\qquad\text{对所有 }w\in T_pM
\end{align*}
定义的余向量。

要看清这为什么是一个光滑丛同态，最容易的办法是考察它在光滑向量场上的作用：
\begin{align*}
\hat g(X)(Y)=g(X,Y)\qquad\text{对 }X,Y\in\mathfrak X(M).
\end{align*}
由于 $\hat g(X)(Y)$ 作为 $Y$ 的函数是 $C^\infty(M)$-线性的，由张量刻画引理（引理\ref{lemma:ch12-tensor-characterization}）可知 $\hat g(X)$ 是一个光滑余向量场；又由于 $\hat g(X)$ 作为 $X$ 的函数是 $C^\infty(M)$-线性的，这就由丛同态刻画引理（引理\ref{lemma:ch10-bundle-hom-characterization}）把 $\hat g$ 定义为一个光滑丛同态。如常，我们对逐点的丛同态 $\hat g:TM\to T^*M$ 和截面上的线性映射 $\hat g:\mathfrak X(M)\to\mathfrak X^*(M)$ 使用同一符号。

注意 $\hat g$ 在每一点都是单射，因为若对某个 $v\in T_pM$ 有 $\hat g(v)=0$，则
\begin{align*}
0=\hat g(v)(v)=\langle v,v\rangle_g,
\end{align*}
而这又蕴含 $v=0$。于是由维数原因，$\hat g$ 是双射，故为一个丛同构（见命题\ref{proposition:ch10-rank-equal}）。

在任意光滑坐标 $(x^i)$ 下，可写 $g=g_{ij}\,\mathrm{d}x^i\mathrm{d}x^j$。于是若 $X,Y$ 是光滑向量场，则
\begin{align*}
\hat g(X)(Y)=g_{ij}X^iY^j,
\end{align*}
它蕴含余向量场 $\hat g(X)$ 具有坐标表达式
\begin{align*}
\hat g(X)=g_{ij}X^i\,\mathrm{d}x^j.
\end{align*}
换言之，$\hat g$ 是这样一个丛同态：它关于 $TM$ 与 $T^*M$ 的坐标标架的矩阵与 $g$ 本身的矩阵相同。（实际上它是 $g$ 的矩阵的转置，但由于 $(g_{ij})$ 对称，二者相同。）

习惯上把余向量场 $\hat g(X)$ 的分量记作
\begin{align*}
\hat g(X)=X_j\,\mathrm{d}x^j,
\qquad\text{其中 }X_j=g_{ij}X^i.
\end{align*}
正因如此，人们说 $\hat g(X)$ 是由 $X$ 通过\emph{降指标}得到的。记号 $X^\flat$ 常用来表示 $\hat g(X)$，因为符号 $\flat$（“降调”）在音乐记号中用来表示把一个音降低。

逆映射 $\hat g^{-1}:T_p^*M\to T_pM$ 的矩阵因而是 $(g_{ij})$ 的逆。（因为 $(g_{ij})$ 是同构 $\hat g$ 的矩阵，它在每一点处可逆。）令 $(g^{ij})$ 表示这样一个矩阵值函数：它在 $p\in M$ 处的值是矩阵 $(g_{ij}(p))$ 的逆，于是
\begin{align*}
g^{ij}g_{jk}=g_{kj}g^{ji}=\delta^i_k.
\end{align*}
由于 $g_{ij}$ 是对称矩阵，$g^{ij}$ 也是，你可以容易地验证。于是对一个余向量场 $\omega\in\mathfrak X^*(M)$，向量场 $\hat g^{-1}(\omega)$ 具有坐标表示
\begin{align*}
\hat g^{-1}(\omega)=\omega^i\frac{\partial}{\partial x^i},
\qquad\text{其中 }\omega^i=g^{ij}\omega_j.
\end{align*}
我们用记号 $\omega^\sharp$（“$\omega$-升调”）表示 $\hat g^{-1}(\omega)$，并说 $\omega^\sharp$ 是由 $\omega$ 通过\emph{升指标}得到的。由于符号 $\flat$ 与 $\sharp$ 借自音乐记号，这两个互逆同构常称为\emph{音乐同构}。一个帮助牢记降调与升调运算的简便记忆法是：$\omega^\sharp$ 在每一点的值是一个向量，我们把它想象成一个（升调的）箭头；而 $X^\flat$ 的值是一个余向量，我们借助它的（降调的）水平集来想象。

升调运算最重要的用途，是把梯度重新确立为 Riemann 流形上的一个向量场。对带边或不带边 Riemann 流形 $(M,g)$ 上任意光滑实值函数 $f$，我们定义一个向量场，称为 $f$ 的\emph{梯度}：
\begin{align*}
\operatorname{grad}f=(\mathrm{d}f)^\sharp=\hat g^{-1}(\mathrm{d}f).
\end{align*}
展开这些定义，可见对任意 $X\in\mathfrak X(M)$，梯度满足
\begin{align*}
\langle\operatorname{grad}f,X\rangle_g=\hat g(\operatorname{grad}f)(X)=\mathrm{d}f(X)=Xf.
\end{align*}
因此 $\operatorname{grad}f$ 是满足
\begin{align*}
\langle\operatorname{grad}f,X\rangle_g=Xf\qquad\text{对每个向量场 }X
\end{align*}
的唯一向量场；等价地，
\begin{align*}
\langle\operatorname{grad}f,\cdot\rangle_g=\mathrm{d}f.
\end{align*}
在光滑坐标下，$\operatorname{grad}f$ 具有表达式
\begin{align*}
\operatorname{grad}f=g^{ij}\frac{\partial f}{\partial x^i}\frac{\partial}{\partial x^j}.
\end{align*}
特别地，这表明 $\operatorname{grad}f$ 是光滑的。在带 Euclid 度量的 $\mathbb{R}^n$ 上，它简化为
\begin{align*}
\operatorname{grad}f=\delta^{ij}\frac{\partial f}{\partial x^i}\frac{\partial}{\partial x^j}
=\sum_{i=1}^n\frac{\partial f}{\partial x^i}\frac{\partial}{\partial x^i}.
\end{align*}
因此在这种情形下我们关于梯度的新定义与初等微积分中的梯度一致。但在其他坐标下，梯度一般不具有相同形式。

\begin{example}\label{example:ch13-gradient-polar}
我们来在极坐标下关于 Euclid 度量计算函数 $f\in C^\infty(\mathbb{R}^2)$ 的梯度。由例\ref{example:ch13-polar-metric}，$\bar g$ 在极坐标下的矩阵是 $\begin{pmatrix}1&0\\0&r^2\end{pmatrix}$，故其逆矩阵是 $\begin{pmatrix}1&0\\0&1/r^2\end{pmatrix}$。把它代入梯度公式，得
\begin{align*}
\operatorname{grad}f=\frac{\partial f}{\partial r}\frac{\partial}{\partial r}+\frac1{r^2}\frac{\partial f}{\partial\theta}\frac{\partial}{\partial\theta}.
\end{align*}
\end{example}

习题\ref{exercise:ch13-21}表明，Riemann 流形上函数 $f$ 的梯度与它在 Euclid 空间中具有相同的几何解释：它的方向是 $f$ 增长最快的方向，并与 $f$ 的水平集正交；它的长度是 $f$ 在任意方向上的最大方向导数。
\end{document}
